\section{Случайные векторы.}
\subsection{Теоретическая справка.}
\begin{definition}
    Случайным вектором из вероятностного пространства $(\Omega, \mathcal{F}, \P)$ называется измеримое отображение $X: \Omega \to \R^n$.
\end{definition}
\begin{note}
    Нетрудно понять, что это определение эквивалентно тому, как если бы мы взяли $n$ случайных величин и составили бы из них отображение.
\end{note}
\begin{definition}
    Функцией распределения случайного вектора $\vec{\xi}$ называется
    \[
        F_{\vec{\xi}}(x_1, \ldots, x_n) = \P(w \in \Omega \mid \xi_1(w) \leq x_1, \ldots, \xi_n(w) \leq x_n) = \P(\xi_1 \leq x_1 \wedge \ldots \wedge \xi_n \leq x_n).
    \]
\end{definition}
\begin{proposition}
    Совместная функция распределения обладает следующими свойствами:
    \begin{enumerate}
        \item $F_{\vec{\xi}}: \R^n \to [0, 1]$.
        \item $F_{\vec{\xi}}$ монотонно неубывающая по каждому аргументу.
        \item $F_{\vec{\xi}}$ непрерывна справа по каждому аргументу.
        \item $F_{\vec{\xi}}(x_1, \ldots, x_n) \ra 1, \|x\| \ra +\infty$.
        \item $F_{\vec{\xi}}(x_1, \ldots, x_n) \ra 0, x_k \ra -\infty$.
        \item $\lim\limits_{x_l \ra +\infty} F_{\vec{\xi}} = F_{\vec{\eta}}$, где $\vec{\eta}$ получен удалением $l$-ой компоненты $\vec{\xi}$.
    \end{enumerate}
\end{proposition}
\begin{definition}
    Математическим ожиданием случайного вектора называется вектор из математических ожиданий его компонент.
\end{definition}
\begin{definition}
    Ковариационной матрицей случайного вектора $\vec{\xi}$ называется матрица $K$, компонентами которой являются $k_{i, j} = \cov(\xi_i, \xi_j)$.
\end{definition}
\begin{proposition}
    Для ковариационной матрицы справедливы следующие свойства:
    \begin{enumerate}
        \item $K^T = K$ и $K$~---~неотрицательно определена.
        \item $\dim K(\R^n) = \rank K$. Значения вектора $\xi$ лежат в $E$ п.н.
    \end{enumerate}
\end{proposition}
\begin{definition}
    Случайные величины $\xi_1, \ldots, \xi_n$ независимы, если их совместная функция распределения равна произведению отдельных функций распределения.
\end{definition}
\begin{note}
    Аналогичное верно и для совместной плотности, если она существует.
\end{note}
\begin{theorem}[О замене меры]
    Пусть $\vec{X}$~---~случайный вектор со значениями в $\R^n$. Пусть $f: \R^n \to \R^1$~---~измерима.
    Тогда, если $f(\vec{X})$~---~интегрируема, то левая и правая части следующего равенства существуют одновременно и совпадают:
    \[
        \Ex f(\vec{X}) = \int_{\R^n} fd\P_{\vec{X}}.
    \]
\end{theorem}
\begin{definition}
    Пусть $g$~---~борелевская функция.
    Функцией $g$ от случайной величины $\xi$ будем называть отображение $\eta = g(\xi) \Lra \eta = g \circ \xi$.
\end{definition}
\begin{theorem}
    Пусть $\xi_1, \ldots, \xi_n$~---~независимые случайные величины, и $\eta_1 = g_1(\xi_1), \ldots, \eta_n = g_n(\xi_n)$~---~случайные величины.
    Тогда $\eta_1, \ldots, \eta_n$~---~независимы.
\end{theorem}
\begin{definition}
    Если для функции $p_{\vec{\xi}}(t_1, \ldots, t_n)$ на $\R^n$ верно
    \[
        F_{\vec{\xi}}(x_1, \ldots, x_n) = \int_{-\infty}^{x_1}\ldots \int_{-\infty}^{x_n}p_{\vec{\xi}}d\vec{x}^n,
    \]
    то $p_{\vec{\xi}}$ называется совместной плотностью распределения.
\end{definition}
\begin{proposition}
    Запишем некоторые свойства совместной плотности:
    \begin{enumerate}
        \item $\P(\vec{\xi} \in A) = \int_A p_{\vec{\xi}}d\vec{x}$ для борелевских множеств $A$.
        \item Если $\xi_1, \ldots, \xi_n$~---~независимые случайные величины, имеющие совместную плотность, то
         \[F_{\vec{\xi}} = F_{\xi_1}\cdot \ldots \cdot F_{\xi_n}\]
        Если взять $n$-ую производную (один раз по каждой компоненте), то с одной стороны мы получим совместную плотность, а с другой -- произведение плотностей.
        Следовательно, совместная плотность независимых случайных величин является произведением их индивидуальных плотностей.
    \end{enumerate}
\end{proposition}
\begin{proposition}
    При гладком диффеоморфизме $\vec{\eta} = \vec{g}(\xi)$ плотность случайного вектора преобразуется как
    \[
        f_{\vec{\eta}}(\vec{y}) = \dfrac{1}{|\det D\vec{g}|}f_{\vec{\xi}}(\vec{g}^{-1}(\vec{y})).
    \]
    В частности, для невырожденного линейного преобразования $A$:
    \[
        f_{\vec{\eta}}(\vec{y}) = \dfrac{1}{|\det A|}f_{\vec{\xi}}(A^{-1}(\vec{y} - \vec{b})).
    \]
\end{proposition}
\begin{proposition}
    Плотность распределения суммы независимых абсолютно-непрерывных случайных величин вычисляется как свёртка плотностей, то есть
    \[
        p_{\xi + \eta}(t) = p_{\xi} * p_{\eta}(t) = \int_{\R} p_{\xi}(x) \cdot p_{\eta}(t - x)dx.
    \]
\end{proposition}
